\scanpage{28}
\mspage{19}
% archive sheet number: 28

% source lines 1--7
 de la même façon qu'un atelier ``spacieux'' (ou ordinaire) se définit en termes d'un magasin quelc. et de $\mathrm{Fig}(M)$. En termes d'un $\unsure{\mathcal F|F}\subset\unsure{\mathcal P|\mathscr P}(M)$, ou plutôt $\mathcal F\subset\mathrm{Fig}_{\unsure{\mathrm{pol}|\mathrm{pré}}}(M)$, les axiomes pertinents, qui \unsure{paraphrasent} Mag at (p.~10) est

% source lines 8--12
\uline{Mag at pol}
\[
\begin{aligned}
\text{a)}\quad &F\in\mathcal F,\ G\leq F\text{ (i.e. }G\text{ partie fermée de}\\
&F\text{ pour }\leq\text{ dans }M)\ \Longrightarrow G\in\mathcal F\\
\text{b)}\quad &F,G\in\mathcal F,\ F\unsure{\gtrless|\lessgtr}_{\mathrm{pol}}G\ \Longrightarrow F\cup G\in\mathcal F\\
\text{c)}\quad &\forall \unsure{x|X}\in M,\ \widetilde x\in\mathcal F\\
\text{c')}\quad &\mathcal F\neq\varnothing\quad(\Longleftrightarrow\varnothing\in\mathcal F).
\end{aligned}
\]

% source lines 13--22
\uline{NB} Un ``atelier spacieux'' \uline{est} \inserted{aussi} un ``atelier polyédral'' ssi le magasin associé est polyédral, \deleted{et si} \deleted{si de plus} (i.e. $X\ll Y\Longrightarrow X\ll_{\mathrm{pol}}Y$) \deleted{et si} (chose qui n'est pas vraie), et si \uline{de plus}, $X\gtrless Y\Longrightarrow X\gtrless_{\mathrm{pol}}Y$ (ce qui est beaucoup plus exceptionnel(?)). Il vaut mieux se faire une raison que ce sont des notions très voisines, mais \unsure{disparates}, et que
% slanted
\margin{Cela provient \unsure{seulement} du point de vue \unsure{choisi}, où on a pris $\ll$ comme relation primitive, mais non $\gtrless$. Il y a là un certain arbitraire\ldots}

% source lines 23--26
 a) Les notions ``spacieuses'' \deleted{\unsure{les}} $\ll$, $\gtrless$ sont plus commodes quand il s'agit de travailler avec des figures aussi générales que possibles, et que

% source lines 27--33
 b) Les notions ``polyédrales'' $\ll_{\mathrm{pol}}$, $\gtrless_{\mathrm{pol}}$ sont plus commodes quand il s'agit de travailler avec des figures, \unsure{vues|vus} comme des ``découpages'' \inserted{\unsure{d'autres|d'icelles}} en ``morceaux'' aussi ``standard'' que possible, avec des ``\unsure{ajustement|ajustements} \deleted{\unsure{standard}} des morceaux'' aussi \deleted{\unsure{S|A}} standard que possible.
